\documentclass[../DoAn.tex]{subfiles}
\begin{document}

\chapter{Nền tảng lý thuyết và Công nghệ sử dụng}
\label{chuong:cong_nghe}

Chương này trình bày các công nghệ, nền tảng cơ sở đã được lựa chọn để giải quyết các bài toán và thiết kế trò chơi được đặt ra ở Chương 2. Với mỗi công nghệ/nền tảng, ưu nhược điểm so với các lựa chọn thay thế sẽ được phân tích để làm rõ lý do sử dụng trong dự án.

\section{Nền tảng phát triển Unreal Engine 5.6}
\label{section:3.1}

Để vận hành một trò chơi 3D với góc nhìn thứ ba, yêu cầu về mặt đồ họa theo thời gian thực (real-time rendering), hiệu năng mô phỏng vật lý và trí tuệ nhân tạo (AI) là vô cùng khắt khe. Dự án sử dụng bộ khung (framework) của Unreal Engine 5.6. Thay vì đi sâu vào hướng dẫn công cụ cấu hình, phần này tập trung phân tích các nền tảng thuật toán cốt lõi bên trong Unreal Engine giải quyết các yêu cầu kỹ thuật của dự án.

\subsection{Chiếu sáng toàn cầu động với Lumen (Lumen Global Illumination)}
Một trong những điểm trọng tâm của đồ án là cơ chế cổng không gian (Portal). Việc kết nối và di chuyển giữa các không gian có điều kiện ánh sáng hoàn toàn khác nhau khiến cho ánh sáng trong màn chơi biến đổi linh hoạt và liên tục.
\begin{itemize}
    \item \textbf{Cơ chế hoạt động:} Lumen là hệ thống Global Illumination (GI) và Reflections động hoàn toàn. Thuật toán của Lumen sử dụng hệ thống \textit{Surface Cache} kết hợp cùng \textit{Screen Tracing} và \textit{Software Ray Tracing} dựa trên Signed Distance Fields (SDF) để nội suy ánh sáng theo thời gian thực.
    \item \textbf{Lựa chọn thay thế:} 
    \begin{itemize}
        \item \textit{Lightmap (Nướng ánh sáng tĩnh):} Phổ biến và tốn ít tài nguyên phần cứng, nhưng không đáp ứng được yêu cầu thay đổi ánh sáng theo thời gian thực khi môi trường bị xáo trộn.
        \item \textit{Hardware Ray Tracing truyền thống:} Nội suy tia sáng thực tế sát với đời thực nhưng yêu cầu tài nguyên thiết bị giới hạn ở các dòng card đồ họa cao cấp. 
    \end{itemize}
    \item \textbf{Lý do hình thành giải pháp:} Lumen giải quyết trọn vẹn bài toán ánh sáng cho hệ thống góc nhìn qua Portal mà vẫn tối ưu chi phí khung hình (framerate) bằng cách giảm thiểu tia trace qua Global Distance Field và tận dụng tài nguyên đệm bộ nhớ (cache).
\end{itemize}

\subsection{Trí tuệ nhân tạo qua Hệ thống Behavior Tree}
Hệ thống kẻ địch (Guard) cần thực hiện hàng loạt hành động xen kẽ như: Tuần tra tự do mốc điểm ngẫu nhiên, Điều tra vùng có âm thanh khả nghi, và Giao tranh khi phát hiện người chơi.
\begin{itemize}
    \item \textbf{Cơ chế hoạt động:} Behavior Tree (Cây Hành vi) phân nhánh việc đánh giá cấu trúc cây ra quyết định từ gốc (Root) qua các nút (Sequence, Selector). Sự kết hợp với dữ liệu chia sẻ (Blackboard) cho phép AI đánh giá môi trường để chọn ra nhánh logic thực thi nhanh nhất.
    \item \textbf{Lựa chọn thay thế - State Machine (Máy trạng thái):} Dễ áp dụng cho AI cơ bản (chỉ có Đứng yên và Tấn công). Tuy nhiên, khi hành vi của AI mở rộng đòi hỏi thêm nhánh Tuần tra, Phỏng đoán (Overwatch), hay Ẩn nấp (Cover), Máy trạng thái có xu hướng phình to không thể kiểm soát thông qua các mạng lưới điều hướng trạng thái chằng chịt (Spaghetti Logic).
    \item \textbf{Lý do hình thành giải pháp:} Cây hành vi với đặc tính mô-đun hóa cao là thuật toán phù hợp nhất. Nhờ đó, dự án dễ dàng thiết kế ba tính cách độc lập cho kẻ thù như Điều tra tích cực (Investigator), Cố thủ (Position Holder), hay Chiến thuật (Tactical) từ các tác vụ (Task) có thể tái sử dụng.
\end{itemize}

\section{Hệ thống kiểm soát phiên bản Perforce Helix Core (P4V)}
\label{section:3.2}

Với kiến trúc trò chơi AAA, toàn bộ tài nguyên (Asset, Texture, Audio) và màn chơi (.umap) thường xuyên ở định dạng nhị phân nguyên khối và nặng vài trăm Megabytes đến hàng Gigabytes. Hệ thống yêu cầu kiểm soát đa người dùng trên các tài nguyên này mà không phá vỡ tính toàn vẹn của mã nguồn.

\subsection{Các nền tảng thay thế}
\begin{itemize}
    \item \textbf{Git (kết hợp Git LFS):} Vô cùng mạnh và linh hoạt ở cấu trúc cây mã nguồn phân tán cho tệp văn bản mã lệnh (C++ source code, Python,...). Tuy nhiên, nếu hai thành viên cùng chỉnh sửa một file bản đồ 3D nhị phân, Git sẽ không thể hợp nhất (merge) kết quả. Việc khắc phục hậu quả hoặc khôi phục tệp trên Git LFS gặp gián đoạn lớn.
    \item \textbf{Subversion (SVN):} Thiết lập quản trị dữ liệu tập trung tuy nhiên tốc độ giao tiếp và gửi tệp giới hạn, năng lực làm việc với hệ sinh thái nhị phân 3D chậm, nhiều rủi ro lỗi xung đột cục bộ.
\end{itemize}

\subsection{Lý do sử dụng Perforce (P4V)}
\begin{itemize}
    \item \textbf{Cơ chế quản trị tập trung và Khóa tệp (File Locking):} Bất kể khi nào một mô hình hoặc màn chơi cần sửa đổi, giao thức Checkout của P4V sẽ lập tức khóa tệp đó lại trạng thái Only-Read (chỉ đọc) đối với người xung quanh. Giải pháp này xóa bỏ 100\% tỷ lệ nhầm lẫn ghi đè tệp nhị phân giữa các thành viên.
    \item \textbf{Tích hợp hệ sinh thái luồng:} Cho phép xử lý và đẩy dữ liệu (Submit) trực tiếp mạnh mẽ qua các tệp Workspace. Khả năng liên kết luồng phát triển dữ liệu của Unreal Engine và C++ hoạt động cực kì hiệu quả và minh bạch cục bộ, hỗ trợ việc tua lại nhanh chóng lịch sử dự án ở mọi quá trình thực nghiệm đồ án.
\end{itemize}


\end{document}